\DocumentMetadata{
  pdfversion=2.0,
  pdfstandard=ua-2,
  lang=en-US,
  tagging=on,
  tagging-setup={math/setup=mathml-SE,tabsorder=structure}
}
\documentclass[11pt,letterpaper]{article}
\usepackage[margin=0.68in,includefoot]{geometry}
\usepackage{newcomputermodern}
\usepackage{amsmath,amscd,mathtools}
\usepackage{tikz,adjustbox}
\providecommand{\square}{\mdlgwhtsquare}
\usepackage[hidelinks]{hyperref}
\hypersetup{
  pdftitle={Algebra First-Year Exam, January 2016, Part 2},
  pdfauthor={University of Florida Department of Mathematics},
  pdfsubject={Graduate mathematics examination},
  pdfdisplaydoctitle=true,
  bookmarksopen=true
}
\setcounter{secnumdepth}{0}
\setlength{\parindent}{0pt}
\setlength{\parskip}{0.28em}
\setlength{\emergencystretch}{3em}
\setlength{\leftmargini}{2.15em}
\setlength{\leftmarginii}{2.65em}
\setlength{\leftmarginiii}{3em}
\renewcommand{\labelenumi}{\textbf{\theenumi.}}
\pagestyle{empty}
\raggedbottom
\allowdisplaybreaks
\newenvironment{examproblems}[1]{%
  \begin{enumerate}%
  \setcounter{enumi}{#1}%
  \setlength{\topsep}{0.35em}%
  \setlength{\partopsep}{0pt}%
  \setlength{\itemsep}{0.48em}%
  \setlength{\parsep}{0.18em}%
}{\end{enumerate}}
\newenvironment{examsubparts}{%
  \begin{enumerate}%
  \setlength{\topsep}{0.18em}%
  \setlength{\partopsep}{0pt}%
  \setlength{\itemsep}{0.16em}%
  \setlength{\parsep}{0.1em}%
}{\end{enumerate}}
\newenvironment{examinstructions}{%
  \par\begingroup\itshape
}{%
  \par\endgroup\vspace{0.18em}
}
\makeatletter
\renewcommand\section{\@startsection{section}{1}{0pt}%
  {-1.25ex plus -.25ex minus -.1ex}%
  {0.55ex plus .1ex}%
  {\normalfont\large\bfseries}}
\renewcommand{\@maketitle}{%
  \newpage\null\vskip -1em%
  {\footnotesize\bfseries
    \href{https://www.ufl.edu/}{University of Florida}, \href{https://math.ufl.edu/}{Department of Mathematics}\par}%
  \vskip -0.5em%
  \tagstructbegin{tag=Title}%
    {\LARGE\bfseries\href{https://gma.math.ufl.edu/exams/algebra-first-year/}{Algebra First-Year Exam}, January 2016, Part 2\par}%
  \tagstructend%
  \vskip 0.25em%
  \tagmcbegin{artifact}%
    \hrule height 1pt%
  \tagmcend%
  \vskip 1em%
}
\makeatother
\title{Algebra First-Year Exam, January 2016, Part 2}
\author{}
\date{}
\begin{document}
\maketitle
\thispagestyle{empty}
\begin{examinstructions}
Answer four problems. (If you turn in more, the first four will be graded.) Within reason, you may quote theorems as long as you state them clearly.
\end{examinstructions}
\begin{examproblems}{0}
\item
Let~\(F\) be a field, let \(n \geq 2\), and let~\(R\) be the ring of \(n \times n\) matrices over~\(F\).
\begin{examsubparts}
\item[\textnormal{(a)}]
Show that the only 2-sided ideals of~\(R\) are \(\{0\}\) and~\(R\).
\item[\textnormal{(b)}]
Show that~\(R\) contains a nonzero proper left ideal~\(I\).
\end{examsubparts}
\item
Give an example of each of the following, with some explanation.
\begin{examsubparts}
\item[\textnormal{(a)}]
A UFD which is not a PID.
\item[\textnormal{(b)}]
A ring~\(R\) with unit and a ring~\(S\) with unit such that~\(S\) is a subring of~\(R\) but \(1_R \neq 1_S\).
\item[\textnormal{(c)}]
A commutative ring~\(R\) with unit and an irreducible element \(\pi \in R\) which is not prime.
\end{examsubparts}
\item
Prove that the ideal \((X,Y)\) in the ring \(\mathbb{Q}[X,Y]\) is not a principal ideal.
\item
Let~\(R\) be a ring and let~\(M\) be an~\(R\)-module. Define 
\[
\operatorname{Tor}(M)=\{x\in M:rx=0\text{ for some }r\in R\setminus\{0\}\}.
\]
\begin{examsubparts}
\item[\textnormal{(a)}]
Prove that if~\(R\) is an integral domain then \(\operatorname{Tor}(M)\) is a submodule of~\(M\).
\item[\textnormal{(b)}]
Give an example of a ring~\(R\) and an~\(R\)-module~\(M\) such that \(\operatorname{Tor}(M)\) is not a submodule of~\(M\). Prove any claims that you make.
\end{examsubparts}
\item
Find a representative for each conjugacy class in \(\operatorname{GL}_2(\mathbb{Z}/3\mathbb{Z})\) consisting of elements whose order divides~\(8\). (Note that we have \(X^4+1=(X^2-X-1)(X^2+X-1)\) in \((\mathbb{Z}/3\mathbb{Z})[X]\).)
\item
Let \(L/K\) be an algebraic field extension and let~\(R\) be a subring of~\(L\) that contains~\(K\). Prove that~\(R\) is a field.
\end{examproblems}
\end{document}
